\documentclass[journal, onecolumn]{IEEEtran}

\usepackage{titlesec}
\usepackage{graphicx}
\usepackage{array}
\usepackage{makecell}
\usepackage{float}
\usepackage{caption}
\usepackage{wrapfig}

\titleformat{\section}
    {\normalfont\bfseries}
    {\thesection}
    {1em}
    {}

\titleformat{\subsection}
    {\normalfont\bfseries}
    {\thesubsection}
    {1em}
    {}

\renewcommand\thesection{\arabic{section}}
\renewcommand\thesubsection{\thesection.\arabic{subsection}}
\renewcommand\thesubsubsection{\thesubsection.\arabic{subsubsection}}

\begin{document}

\title{FireFly Design Document}
\author{Christian Li, Kenny Wang, Kevin Yao, Matin Sobhani, Nick Lafredo, Parsa Rezvani, Ryan Yu}
\date{\today}

\maketitle

\begin{figure*}[!t]
    \centering
    \includegraphics[width=.5\linewidth]{images/UCSC_Logo.png}
    \caption*{}
    \label{fig:UCSCLOGO}
\end{figure*}

\vspace{1in}

\begin{figure*}[!t]
    \centering
    \includegraphics[width=0.25\linewidth]{images/BE_Logo.png}
    \caption*{}
    \label{fig:BELOGO}
\end{figure*}

\clearpage

\tableofcontents

\vspace{1in}

\begin{abstract}
This document explains how to design, build, and operate the FireFly FireWatch Drone.
\end{abstract}

\begin{IEEEkeywords}
none
\end{IEEEkeywords}

\newpage

% ============================================================
\section*{Executive Summary}
% ============================================================

FireFly is a rapidly deployable quadcopter platform designed as a first-response environmental monitoring tool for high wildfire-risk areas with little or no sensor infrastructure. The system integrates temperature, humidity, and smoke sensing with live data transmission to a ground station, enabling early awareness and alerts when fire-risk is detected. Its mechanical and electrical architecture allows it to adapt to different terrains and sensing requirements.

The design section of the document specifies system requirements for payload capacity, flight time, communication range, and sensor capability. It targets operation in terrain suited poorly for manual monitoring. The architecture defines the flight controller, power distribution, propulsion stack, communication link, and sensor bus. The mechanical design focuses on a 3D-printed frame with mounting points for motors, batteries, and payloads.

Power and propulsion modules are chosen to maximize flight time while keeping a stable rate of power consumption. Sensor and payload modules include environmental sensors that measure temperature, humidity, and smoke concentration to support fire monitoring. The electronics and control subsection describes how the flight controller interfaces with ESC, receivers, and sensors over buses.

Safety and reliability are assured through design features and planned validation. Propeller guards and keyed connectors prevent user error and physical hazards. The document outlines bench testing for individual modules (power, propulsion, sensors, communications) and flight testing for stable hover, controlled movement, data link, and detection of smoke and temperature. These elements aim to demonstrate that FireFly can serve as a practical platform for wildfire monitoring and emergency response research.

\clearpage

% ============================================================
\section{Introduction}
% ============================================================

This document contains all the information needed to design and use the FireFly Drone.

\vspace{1em}

Why is collecting data during a pre-fire window important? A pre-fire window---the period before an incident occurs---is critical because it provides time-sensitive information that improves responder safety, reduces property damage, and enables faster, more effective emergency response. Data on vegetation conditions, topography, temperature, humidity, and smoke concentration helps predict fire behavior, identify high-risk areas, and support mitigation actions such as controlled burns or targeted thinning operations. Without this information, emergency crews arrive with incomplete situational awareness and must gather critical details in real time under dangerous conditions.

\vspace{1em}

For wildfire response specifically, the challenge is compounded by geography. High-risk areas are frequently roadless wilderness where permanent sensor infrastructure cannot be installed or maintained. Satellite imagery and aerial patrols provide coverage, but both suffer from response delays and gaps during the highest-risk windows. The result is a critical blind spot: the pre-fire period in remote terrain is precisely where early awareness matters most, yet it is the hardest environment in which to collect data.

\vspace{1em}

FireFly addresses this gap directly. By deploying a sensor-equipped quadcopter into areas with zero existing infrastructure, FireFly enables real-time collection of temperature, humidity, and smoke data during the pre-fire window, transmitting live readings to a ground station and triggering alerts when fire risk thresholds are crossed. The platform is designed for rapid deployment by field crews without requiring extensive training or specialized support, transforming what would otherwise be a period of uncertainty into a structured, data-driven response.

\subsection{Need Statement}
Remote and high fire-risk areas lack the sensor infrastructure necessary to monitor environmental conditions during the critical pre-fire window, leaving fire response teams without the real-time data needed to act early.

\subsection{Goal Statement}
Provide fire response teams with real-time temperature, humidity, and smoke readings from high-risk, infrastructure-free terrain by deploying a rapidly mobilized monitoring system that transmits live data to a ground station and alerts operators when fire risk thresholds are crossed.

\subsection{Personas}
Individuals with experience levels ranging from novice to expert in fire monitoring and environmental fieldwork.

\subsubsection*{End Users}

\begin{wrapfigure}{l}{0.2\textwidth}
    \includegraphics[width=0.15\textwidth]{images/marcus.png}
\end{wrapfigure}
\textbf{Marcus} \\
\textit{Role:} Wildfire Operations Coordinator (professional first responder) \\
\textit{Background / Context:} Works with a state forestry division managing prescribed burns and post-fire monitoring. Needs something fast that he can hand off to field crews without extensive training. \\
\textit{Goals:} Rapid deployment of sensor-equipped drones into areas too dangerous or inaccessible for personnel during the high-risk pre-fire window. Wants real-time smoke and temperature data without putting people in harm's way. \\
\textit{Behaviors:} Works under time pressure, values reliability over features, and coordinates with multiple agencies. Skeptical of unproven technology. Practical --- if it breaks in the field, it needs to be fixable on-site with minimal tools. \\
\textit{Expectations:} A rugged, simple-to-operate platform that can be deployed at a moment's notice. Wants clear threshold alerts he can act on immediately.

\vspace{1em}

\begin{wrapfigure}{l}{0.2\textwidth}
    \includegraphics[width=0.18\textwidth]{images/elena.png}
\end{wrapfigure}
\textbf{Elena} \\
\textit{Role:} Ecology Graduate Researcher (academic user) \\
\textit{Background / Context:} Studying fire risk and ecosystem conditions at a land-grant university. Needs to collect environmental sensor data across wilderness areas where no permanent monitoring infrastructure exists. Currently hikes in with handheld sensors, which limits how often and how broadly she can sample. \\
\textit{Goals:} Automate pre-fire data collection over defined survey areas. Wants live temperature, humidity, and smoke readings logged without requiring physical presence at each site. \\
\textit{Behaviors:} Comfortable with technical documentation, participates in academic drone and environmental monitoring communities, and documents everything for publication. Willing to configure and customize but does not have time to build from scratch. \\
\textit{Expectations:} A reliable platform with a clear data pipeline she can connect to her existing research database, and results trustworthy enough to cite.

\vspace{1em}

\subsubsection*{Experts / Consultants}

\begin{wrapfigure}{l}{0.2\textwidth}
    \includegraphics[width=0.18\textwidth]{images/anita.png}
\end{wrapfigure}
\textbf{Dr. Anita Rowe} \\
\textit{Role:} Environmental Sensor Systems Engineer \\
\textit{Background / Context:} Works at a national lab developing atmospheric sensor arrays for wildfire monitoring. Has deep expertise in what data is actually actionable in pre-fire environments and where sensors fail under heat, humidity, and particulate exposure. \\
\textit{Goals:} Validate that FireFly's sensor suite captures meaningful data for fire risk assessment and identify unreliable component choices before they are locked into the design. \\
\textit{Behaviors:} Data-driven and skeptical of consumer-grade sensors in harsh field conditions. Will probe calibration drift, operating temperature limits, and whether threshold alert logic reflects real-world fire risk indicators. \\
\textit{Expectations:} A platform sound enough in its sensor selection and alert logic that field results can be trusted and the methodology can be defended to other scientists.

\vspace{1em}

\subsubsection*{Client}

\begin{wrapfigure}{l}{0.2\textwidth}
    \includegraphics[width=0.18\textwidth]{images/usfs.png}
\end{wrapfigure}
\textbf{U.S. Forest Service Regional Office} \\
\textit{Role:} Federal Land and Fire Management Agency \\
\textit{Background / Context:} Responsible for fire prevention, monitoring, and response across millions of acres of national forest, much of it roadless wilderness with no existing sensor infrastructure. Currently relies on satellite imagery, fire lookout towers, and aerial patrols, all of which have coverage gaps and response delays during the highest-risk windows. \\
\textit{Goals:} Evaluate FireFly as a rapidly deployable supplement to existing monitoring systems for areas where infrastructure cannot operate. Interested in whether the platform can reliably trigger alerts before a fire establishes. \\
\textit{Behaviors:} Evaluates tools based on operational reliability and ease of use for non-specialist field crews. Requires demonstrated performance data before recommending adoption. \\
\textit{Expectations:} A working prototype that proves the core concept --- zero-infrastructure deployment, live sensor transmission, and threshold-based alerting --- with documentation clear enough for field personnel to operate without engineering support.

\subsection{Existing Research}
Drone-based wildfire monitoring is an active research area, yet a clear gap remains for rapidly deployable, infrastructure-independent platforms targeting the pre-fire window in remote terrain.

Current detection methods rely on satellite imagery, fixed camera-tower networks, and manned aircraft. Satellites cannot reliably detect fires smaller than approximately 18 square kilometers, which limits early-stage intervention. Fixed towers such as AlertCalifornia's ~1,000-camera network are constrained by line of sight, leaving roadless and mountainous terrain uncovered. Research shows UAVs can bridge these gaps; for example, studies have demonstrated UAV--IoT networks detecting fires as small as 2.5 square kilometers with greater than 99 percent accuracy. However, many existing platforms depend on pre-deployed ground sensors, persistent connectivity infrastructure, or large multi-drone fleets, creating a blind spot in remote areas where such infrastructure is impractical.

FireFly targets the specific gap the literature acknowledges but has not fully addressed: a self-contained quadcopter requiring no pre-existing infrastructure, no persistent network, and no specialized operator training to deliver real-time environmental sensor data and threshold-based fire risk alerts.

\subsection{Sustainability Statement}

The FireFly project is committed to the ``Design for Environment'' philosophy. Our goal is to provide a high-tech solution for wildfire mitigation while ensuring that the lifecycle of the drone itself does not offset its environmental benefits.

\vspace{1em}

\textbf{Environmental Stewardship:}
The primary purpose of FireFly is the preservation of biodiversity and carbon sinks by preventing catastrophic wildfires. By providing early-warning data in the pre-fire window, we reduce the need for large-scale, carbon-intensive emergency responses such as heavy tanker aircraft and ground fleets. Furthermore, by utilizing electric propulsion over internal combustion, we ensure zero-emission operation during monitoring flights.

\vspace{1em}

\textbf{Material Selection and Circularity:}
To minimize our manufacturing footprint, we have opted for 3D-printed components using \textbf{PLA}. PLA is a bio-based, industrially compostable polymer derived from renewable resources like corn starch. By using additive manufacturing, we significantly reduce material waste compared to traditional subtractive machining.

\vspace{1em}

\textbf{Life Cycle and Battery Management:}
We recognize that Lithium-Polymer (LiPo) batteries pose a chemical risk if disposed of improperly. Our operational plan includes a strict end-of-life protocol where degraded cells are discharged and taken to certified e-waste recycling centers. We have designed the electronics with standardized connectors to ensure that sensors and flight controllers can be harvested and reused in future iterations of the platform, extending the functional life of the high-value silicon components.

\vspace{1em}

\textbf{Social and Safety Impact:}
Sustainability includes the protection of human capital. FireFly enhances the safety of first responders by providing situational awareness in hazardous terrains without risking human life. By creating a rapidly deployable tool, we democratize access to advanced fire-monitoring technology for smaller forestry divisions and academic researchers who may lack the budget for satellite-based or manned aerial surveillance.

\newpage

% ============================================================
\section{Design}
% ============================================================

\subsection{Aesthetic Prototype}

The FireFly aesthetic prototype follows a field-ready design language: a compact enclosed central body with rounded corners, symmetric X-arm geometry, and cylindrical motor pods at each corner. The overall form emphasizes structural clarity and quick visual inspection rather than cosmetic styling. Rounded exterior edges reduce snag points during transport, while the open-top bay and side windows communicate where electronics, wiring, and battery hardware are intended to be serviced.

\begin{figure}[H]
    \centering
    \includegraphics[width=0.46\textwidth]{images/drone_design_5_top.png}
    \hfill
    \includegraphics[width=0.46\textwidth]{images/drone_design_5_corner.png}
    \caption{FireFly CAD aesthetic prototype, iteration v5 (top and isometric views).}
    \label{fig:drone_v5_aesthetic}
\end{figure}

The prototype evolved through multiple print-test iterations. Early versions validated basic packaging and motor placement, while later versions introduced a cleaner body-to-arm transition, larger internal volume, and improved vent/window geometry for lighter weight and faster access to internals. This iterative process preserved a consistent visual identity while converging on a more manufacturable and maintainable frame.

\begin{figure}[H]
    \centering
    \includegraphics[width=0.85\textwidth]{images/drone_design_5_compare_1.jpeg}
    \caption{Physical prototype progression (v1--v3) that informed the final v5 frame direction.}
    \label{fig:drone_proto_progression}
\end{figure}

\subsection{Design for Manufacture, Assembly, Maintenance}

\textbf{Design for Manufacture (DFM):}
The current frame is designed as a single-piece 3D print to minimize part count and remove alignment errors. The v5 CAD envelope is approximately $100 \times 72 \times 24$ mm with $2.5$ mm floor and wall thicknesses, balancing print speed, stiffness, and mass. Arms are sized at $10.0$ mm width and $5.4$ mm thickness to maintain rigidity under motor thrust and landing loads. Motor pods are dimensioned for SunnySky R1106-class motors with built-in tolerance ($19$ mm outer diameter, $15.0$ mm inner diameter, $17$ mm height, $2.2$ mm mounting holes on a $9$ mm bolt-circle).

\vspace{0.5em}

\textbf{Design for Assembly (DFA):}
Assembly features are integrated directly into the print geometry. The frame includes four battery strap slots in the floor and continuous wire-routing paths from motor pods into the electronics bay. This allows a straightforward assembly order: install motors in pods, route motor leads through arm channels, mount the controller board, then secure the battery and sensor harness. Because most interfaces are geometric rather than adhesive, rework is possible without damaging the frame.

\vspace{0.5em}

\textbf{Design for Maintenance (DFMn):}
Maintenance access is provided by the open top cavity, side-wall windows, and pod-to-arm wire notches. These features allow quick inspection of solder joints, connector seating, and wire wear without full teardown. Electronics can be removed independently from propulsion hardware, and battery straps can be replaced in the field with standard hook-and-loop straps. In case of structural damage, the one-piece architecture supports rapid replacement by reprinting the frame using the same CAD file, preserving dimensional repeatability across repairs.

\subsection{Block Diagrams}

\begin{figure}[htbp]
    \centering
    \includegraphics[width=0.5\textwidth]{images/Schematic_Design.png}
    \caption{A high level schematic of our drone.}
    \label{fig:schematic_design}
\end{figure}

\subsection{Wiring Diagrams}

\begin{figure}[htbp]
    \centering
    \includegraphics[width=0.75\textwidth]{images/Circuit Diagram.png}
    \caption{Circuit Diagram}
    \label{fig:circuit_diagram}
\end{figure}

\vspace{1in}

\subsection{State Transition Diagrams}

\begin{figure}[H]
    \centering
    \includegraphics[width=0.75\textwidth]{images/State Transition Diagram.png}
    \caption{State Transition Diagram}
    \label{fig:state_transition}
\end{figure}

\subsection{Technology (Languages, Hardware, etc.)}

FireFly uses a two-node ESP32-C3 architecture connected via the ESP-NOW protocol. The drone carries an ESP32-C3 microcontroller that reads sensor data from a BME280 (temperature and humidity) and an MQ-2 (smoke detection), packages the readings into a structured telemetry packet, and transmits it wirelessly to a ground receiver using ESP-NOW. ESP-NOW is a connectionless peer-to-peer protocol operating on the 2.4GHz band with a range of approximately 400 meters line-of-sight and latency under 10ms, making it well suited for remote wildfire monitoring scenarios where reliable low-latency telemetry is critical.

\vspace{0.5em}

The ground station consists of a second ESP32-C3 that registers as an ESP-NOW receiver, unpacks the incoming telemetry struct, and forwards the data as a JSON stream over USB serial to a host machine. The host machine runs a Python Flask server that reads the serial stream and exposes a local HTTP endpoint, which a browser-based dashboard polls at one-second intervals to display live sensor readings. Firmware on both ESP32-C3 devices is written in C using the ESP-IDF framework, with the ground station dashboard implemented in Python and JavaScript.

\subsection{Simulations}

Before building and flying, we used simulation tools to verify that the FireFly configuration can sustain flight, carry its sensor payload, and meet the core mission requirement of delivering real-time environmental data during the pre-fire window. An online multicopter calculator (eCalc xcopterCalc) was used to evaluate combinations of frame size, motors, propellers, and battery against thrust-to-weight ratio, hover throttle, current draw, and expected flight time under sensor payload assumptions. This confirmed that the selected hardware provides sufficient thrust margin to carry the BME280 temperature/humidity sensor and MQ-2 smoke detector while maintaining acceptable hover efficiency and a flight time long enough to collect and transmit meaningful readings to the ground station.

\vspace{0.5em}

A simplified quadcopter dynamics model was used to tune attitude and altitude control parameters, verifying that the platform can hold a stable position over a monitoring area --- a prerequisite for consistent sensor sampling. Research on quadcopter modeling confirms that such simulations can predict controller behavior and reduce the number of risky test flights needed to achieve stable flight. Together, these simulations provide evidence that FireFly can meet its need statement requirements: reaching areas with zero sensor infrastructure and sustaining flight long enough to collect and transmit temperature, humidity, and smoke data before committing to full implementation and field testing.

\newpage

% ============================================================
\section{Parts List}
% ============================================================

\begin{table}[H]
    \centering
    \small
    \begin{tabular}{|l|c|c|}
        \hline
        \textbf{Part} & \textbf{Supplier} & \textbf{Quantity} \\
        \hline
        CarbonX PLA+CF Filament & Creality & 1 \\
        \hline
        Battery: LiPo 850 mAh 3S 25/30C & HobbyKing & 1 \\
        \hline
        Motor: SunnySky R1106 FPV Brushless Motors & Druav & 4 \\
        \hline
        Propeller: Gemfan Hurricane 2009-3 Tri-Blade (Set of 4) & Pyrodrone & 1 \\
        \hline
        Caddx Ant Lite Analog Camera (FPV Cycle Edition) & CaddxFPV & 1 \\
        \hline
        TFMini-S Micro LiDAR Module & SparkFun & 1 \\
        \hline
        HAKRC 20A 4-in-1 \verb|3-6S BLHeli_32| 5V 3A BEC DShot1200 Mini ESC & Alibaba & 1 \\
        \hline
        Perforated Circuit Board & Adafruit & 2 \\
        \hline
        Female Header Pins 2x6 & AliExpress & 4 \\
        \hline
        Female Header Pins 2x8 & AliExpress & 2 \\
        \hline
        Female Header Pins 2x10 & Mikroe & 2 \\
        \hline
    \end{tabular}
    \caption{Itemized Budget}
    \label{tab:budget}
\end{table}

\newpage

% ============================================================
\section{Appendix 1 - Problem Formulation}
% ============================================================

\subsection{Conceptualisations}

The conceptualization of FireFly establishes the high level structure of the system before detailed design and component selection. FireFly is envisioned as a lightweight, rapidly deployable quadcopter that operates as a mobile environmental sensing node in areas with no existing infrastructure. At this stage, the system is defined in terms of its required capabilities rather than specific hardware.

\vspace{0.5em}

The platform is conceptualized around three core functions: sustained stable flight, reliable environmental sensing, and real time data transmission. The airframe provides mounting points for a modular payload consisting of the temperature, humidity, and smoke sensors. A central flight controller manages propulsion, stabilization, and stability controller communication, while the microcontroller subsystem handles sensor acquisition and mesh network data transport. Power is supplied by a LiPo battery sized to support both propulaion and the onboard electronics for the duration of a pre-fire monitoring mission. 

\vspace{0.5em}

To support use in remote terrain, the communication layer is conceptualized as a decentralized Wi-Fi mesh network rather than a point to point link. This allows each drone to ace as both a sensor playform and a relay, ensuring data can reach the ground station even when terrain blocks direct line of sight. The software model includes a lightweight telemetry server on each node and a centralized logging system at the base station.

\vspace{0.5em}

Overall, the conceptual model defines FireFly as an infrastructure indepenedant quadcopter platform that combines environmental sensing, onboard processing, and long range mesh communication. This provides the foundation for selecting hardware, designing the frame, and implementing the control architecture.

\subsection{Brainstorming}

The brainstorming phase explored multiple design directions before converging on the current FireFly configuration. For the frame, several geometries were evaluated, including traditional X and H-style quadcopter layouts. The final frame concept emerged as a hybrid, leaning primarily toward an X-frame for improved stability, balanced motor distribution, and ease of maintenance while still allowing flexible mounting options for sensors and electronics.

\vspace{0.5em}

The sensing strategy originated from studying existing ground-based mesh networks deployed in forests for environmental monitoring. These networks offer valuable data but are geographically limited and expensive to expand. This led to the idea of using a drone as a mobile extension of such networks, particularly in areas where no sensor infrastructure exists. While commercial sensor networks rely on stationary nodes, FireFly would act as a mobile node capable of collecting and relaying data beyond the reach of fixed installations.

\vspace{0.5em}

Several communication methods were explored. Bluetooth Low Energy (BLE) was initially attractive due to built in characteristics support and team familiarity. Wi-Fi was evaluated next and ultimately selected because it integrates cleanly with the controller app, supports higher throughput, and enables a mesh topology through ESP-Now and Wi-Fi mesh extensions. This allowed the team to conceptualize a swarm capable of relaying data across larger distancees where a single radio link would fail.

\vspace{0.5em}

Power and propulsion options were compared through a combination of bench testing and simulation. Different Lipo battery configurations (3-cell vs. 4-cell), motor-propeller pairs, and ESC compatibility with the microcontroller were evaluated. Using tools such as eCalc, the team analyzed thrust to weight ratio, hover current, and projected flight time to determine whether each candidate configuration could sustain flight with the expected payload. These explorations guided the selection of a power system that balances endurance, stability, and the electrical needs of the sensor and communication hardware.

\vspace{0.5em}

The brainstorming process defined the feasible design space and informed the selection of the frame, mobile sensing concept, Wi-Fi/ESP-Now mesh communication, and optimized power system used in the FireFly design.

\subsection{Decision Tables}

\begin{table}[H]
    \centering
    \begin{tabular}{|l|c|c|c|}
        \hline
        \textbf{Criteria} & \textbf{Option A (High KV)} & \textbf{Option B (Low KV)} & \textbf{Weight} \\
        \hline
        Thrust-to-Weight & 5 & 3 & 40\% \\
        \hline
        Efficiency & 5 & 5 & 40\% \\
        \hline
        Air Time & 5 & 8 & 20\% \\
        \hline
        \textbf{Total Score} & \textbf{4.0} & \textbf{4.0} & \textbf{100\%} \\
        \hline
    \end{tabular}
    \caption{Motor Selection Decision Matrix}
\end{table}

\subsection{Morphological Charts}
The following combinations were explored for the primary subsystems:
\begin{itemize}
    \item \textbf{Frame:} ...
    \item \textbf{Telemetry:} ...
\end{itemize}

\newpage

% ============================================================
\section{Appendix 2 - Planning}
% ============================================================

\subsection{Basic Plan / Gantt Chart}
[...]

\subsection{Division of Labor}

\begin{table}[H]
    \centering
    \begin{tabular}{|l|l|}
        \hline
        \textbf{Team Member} & \textbf{Primary Responsibility} \\
        \hline
        Christian Li & Flight Controller Programming \\
        \hline
        Kenny Wang & Hardware Integration \\
        \hline
        Kevin Yao & Mechanical Design / CAD \\
        \hline
        Matin Sobhani & Mechanical Design / CAD \\
        \hline
        Nick Lafredo & Power Systems / Documentation \\
        \hline
        Parsa Rezvani & Testing and Validation / Product Manager \\
        \hline
        Ryan Yu & Telemetry and Documentation \\
        \hline
    \end{tabular}
    \caption{Division of Labor during Prototyping}
\end{table}

\subsection{Collaboration}
[...]

\newpage

% ============================================================
\section{Appendix 3 - Test Plan \& Results}
% ============================================================

The test plan consists of three phases:
\begin{enumerate}
    \item \textbf{Bench Testing:} ...
    \item \textbf{Tethered Flight:} ...
    \item \textbf{Open Field Flight:} ...
\end{enumerate}

\newpage

% ============================================================
\section{Appendix 4 - Review}
% ============================================================

\textit{Reflection of each team member's thoughts on the project.}

\vspace{1em}

\textbf{Christian Li:} Looking back on this project, one of the biggest lessons I learned is the importance of starting with a clearly defined need rather than a preconceived solution. We approached the project wanting to build a drone, and only later did we step back and ask why a drone was the right solution for the problem. If we had focused earlier on the lack of fire risk data in remote areas, we could have identified the required characteristics sooner and avoided some of the design and task adjustments we made midway through the quarter. Even so, our team communicated well and adapted quickly. The flexibility helped us stay on track and produce a coherent design. This experience also gave me valuable insight into working as a team, improving how we communicate, distribute work, and learn new tools and environments together.

\vspace{1em}

\textbf{Kenny Wang:} Restarting this project I think that we should focus on not building a product that interests us but develop the idea based on a universal problem or a need. From there we can then focus on implementing it in a way that can handle this universal problem or need. For example, if we had started with the problem being that we have no way of reaching remote areas with no fire detection, we can then build upon that and not just specifically building a flying drone but maybe a rover, or even something that does not even have to move. There would be a variety of implementations of helping alleviate the same problem and need while having numerous amounts of angles to implement that solution. Even with some setbacks that we faced during the beginning stages of the project we were able to have good communication within the team to implement these changes so that we were able to complete our design work.

\vspace{1em}

\textbf{Kevin Yao:} [Reflection text here...]

\vspace{1em}

\textbf{Matin Sobhani:} If we were starting this project again, I would make sure we had a solid need/goal statement, and verified it with the TA and professor. We thought we had a solid one early on around week 4, but ended up having to change it around week 8. Thankfully this didn't affect us negatively too much considering our underlying prototype idea is still the same. What went well is the fact that we were very communicative throughout the project and always kept up-to-date.

\vspace{1em}

\textbf{Nick Lafredo:} [Reflection text here...]

\vspace{1em}

\textbf{Parsa Rezvani:} This project has been a significant learning experience for both me and the team. Initially, we approached it as an exercise in selecting an interesting idea and producing documentation around it. Over time, however, we came to understand that the class demanded something closer to real world engineering practice; Defining a defensible need, building toward a clear goal, and maintaining a concept that could stand on its own merits in a professional context. Had we internalized that framing earlier, our process would have been more focused from the start. Despite the adjustments along the way, I am proud of how the team collaborated and adapted. With our components arriving soon, I am confident we are well positioned to make meaningful progress on the prototype.

\vspace{1em}

\textbf{Ryan Yu:} [Reflection text here...]

\end{document}
